\section{Threat Model}
\label{sec:threat-model}

\subsection{Adversary}

We consider an adversary who can publish a package version to a registry. This
capability is obtained in three ways that we observe in the wild: compromising
an existing maintainer account, typically through credential stuffing or a
recovery-email takeover~\citep{mbeki2022typosquat}; acquiring maintainership of
an abandoned package, which several registries grant on request; and
registering a new name that a victim will plausibly install, whether by
typosquatting or by dependency confusion against an internal
name~\citep{ferrante2021incidents}.

The adversary does not control the registry, the network path, or the victim's
machine before the install. They do not need to. The install client fetches
their archive over an authenticated channel and then executes their hook,
which is the entire attack.

\subsection{Victim and assets}

The victim is any machine that installs the package: a developer workstation, a
continuous integration runner, or a container build. Continuous integration is
the most valuable target and the least defended, because a runner typically
holds deployment credentials, registry publishing tokens, and cloud role
credentials reachable from an unauthenticated metadata endpoint, and because
nobody is watching the terminal.

The assets we instrument for are therefore credential material in the
environment and on disk, the cloud metadata endpoint, the source tree being
built, and outbound network reachability. Figure~\ref{fig:threat} shows the
positions and the trust boundaries.

\begin{figure}[t]
  \centering
  \begin{tikzpicture}[
      node distance=5mm and 6mm,
      blk/.style={draw,rounded corners=1.5pt,align=center,font=\scriptsize,
                  minimum height=7mm,minimum width=16mm},
      adv/.style={blk,fill=black!12,draw=black},
      flow/.style={-{Stealth[length=1.7mm]},semithick},
      atk/.style={-{Stealth[length=1.7mm]},semithick,dashed},
    ]
    \node[adv] (att) {Adversary};
    \node[blk,right=12mm of att] (acct) {Maintainer\\account};
    \node[blk,right=13mm of acct] (reg) {Registry};
    \node[blk,below=9mm of reg] (client) {Install client};
    \node[blk,left=of client] (hook) {Hook process\\\scriptsize user privileges};
    \node[blk,below=9mm of hook] (creds) {CI secrets,\\tokens, metadata};
    \node[adv,left=17mm of hook] (drop) {Collection\\host};

    \draw[atk] (att) -- node[above,font=\scriptsize] {takeover} (acct);
    \draw[flow] (acct) -- node[above,font=\scriptsize] {publish} (reg);
    \draw[flow] (reg) -- node[right,font=\scriptsize] {fetch} (client);
    \draw[flow] (client) -- node[above,font=\scriptsize] {exec} (hook);
    \draw[flow] (hook) -- node[right,font=\scriptsize] {read} (creds);
    \draw[atk] (hook) -- node[above,font=\scriptsize] {exfiltrate} (drop);
    \draw[atk] (drop) to[bend left=38] (att);

    \begin{scope}[on background layer]
      \node[draw,dashed,rounded corners=2pt,fit=(client)(hook)(creds),
            inner sep=3mm,
            label={[font=\scriptsize]below:victim build host}] {};
      \node[draw,dotted,rounded corners=2pt,fit=(acct)(reg),
            inner sep=2.5mm,
            label={[font=\scriptsize]above:registry trust boundary}] {};
    \end{scope}
  \end{tikzpicture}
  \caption{Threat model. The adversary never touches the network path or the
  victim host directly. The install client fetches an authenticated archive and
  then executes the adversary's hook with the invoking user's privileges,
  inside the victim's trust boundary.}
  \label{fig:threat}
\end{figure}

\subsection{Out of scope}

We do not consider registry compromise, transport interception, or malicious
behaviour in the package's runtime code. The last is a genuine and larger
problem, but it requires the victim to call the code and is therefore reachable
by ordinary code review, whereas a hook executes before anyone has read
anything. We also exclude hooks that require interactive input, 0.3 percent of
our corpus, since our sandbox cannot supply it.
